\chapter{System Workflow}
\label{chap:workflow}

\section{Workflow Overview}
The DNA Tag and Trace workflow consists of five clear operational stages (Figure~\ref{fig:workflow_diagram}).

\begin{figure}[H]
    \centering
    \begin{tikzpicture}[
        scale=0.9, transform shape,
        stepnode/.style={rectangle, rounded corners=2pt, draw=black, fill=white, line width=0.8pt, text width=6.2cm, minimum height=1.1cm, align=left, font=\small},
        arrowstyle/.style={-{Stealth[length=2.5mm]}, line width=0.9pt, color=black}
    ]
        % Left Column
        \node[stepnode] (s1) at (0, 3.4) {\textbf{1. DNA Designing}\\ \footnotesize Algorithmic generation, $T_m$ matching, BLAST};
        \node[stepnode] (s2) at (0, 1.7) {\textbf{2. Physical Synthesis}\\ \footnotesize PCR scaling, 1L carrier buffer formulation};
        \node[stepnode] (s3) at (0, 0.0) {\textbf{3. Product Integration}\\ \footnotesize Spraying on packaging, liquids \& security stamps};
        
        % Right Column
        \node[stepnode] (s4) at (7.4, 3.4) {\textbf{4. Field Recovery}\\ \footnotesize Surface swabbing \& chemical recovery};
        \node[stepnode] (s5) at (7.4, 1.7) {\textbf{5. qPCR Verification}\\ \footnotesize 4-channel multiplex detection \& deconvolution};
        \node[stepnode] (s6) at (7.4, 0.0) {\textbf{6. Database Match}\\ \footnotesize Ledger query \& verification report};
        
        % Sequential Flow
        \draw[arrowstyle] (s1) -- (s2);
        \draw[arrowstyle] (s2) -- (s3);
        \draw[arrowstyle] (s3.east) -- ++(0.6,0) |- (s4.west);
        \draw[arrowstyle] (s4) -- (s5);
        \draw[arrowstyle] (s5) -- (s6);
    \end{tikzpicture}
    \caption{The five stages of the DNA tagging and verification lifecycle.}
    \label{fig:workflow_diagram}
\end{figure}

\section{DNA Designing}
DNA constructs are designed through the custom software engine:
\begin{itemize}
    \item \textbf{Custom Sizing}: The software supports arbitrary, user-defined sequence lengths according to industrial application requirements (for this proof-of-concept validation, a 200~bp prototype construct was engineered and synthesized).
    \item \textbf{Structure}: Flanked by 20~bp forward and reverse primers, containing four internal 24~bp fluorescent probe targets (FAM, HEX, Texas Red, Cy5) separated by spacers (Figure~\ref{fig:molecular_anatomy}).
    \item \textbf{Thermodynamic Boundaries}: Melting temperatures ($T_m$) are calculated using SantaLucia parameters:
    \begin{equation}
        T_m = \frac{\Delta H^\circ}{\Delta S^\circ + R \ln(C_T / 4)} + 16.6 \log_{10}[\text{Na}^+]
    \end{equation}
    Primers are matched at $T_m = 58.5 \pm 1.0^\circ\text{C}$, and probes are set at $T_m \ge 68.0^\circ\text{C}$ ($+10.0^\circ\text{C}$ above primers).
    \item \textbf{Safety Screening}: Local BLAST checks ensure zero homology ($<25\%$) with any known environmental or human pathogens.
\end{itemize}

\begin{figure}[H]
    \centering
    \begin{tikzpicture}[
        scale=0.88, transform shape,
        boxstyle/.style={
            rectangle, 
            draw=black, 
            line width=0.8pt, 
            minimum width=2.15cm, 
            text width=2.05cm, 
            minimum height=1.35cm, 
            align=center, 
            font=\footnotesize\bfseries
        },
        primerbox/.style={boxstyle, fill=gray!15},
        probebox/.style={boxstyle, fill=white}
    ]
        \node[primerbox] (fwd) at (0, 0) {5' Forward\\ Primer\\ (20 bp)};
        \node[probebox, right=0.12cm of fwd] (p1) {Probe 1\\ FAM\\ (24 bp)};
        \node[probebox, right=0.12cm of p1] (p2) {Probe 2\\ HEX\\ (24 bp)};
        \node[probebox, right=0.12cm of p2] (p3) {Probe 3\\ TxRed\\ (24 bp)};
        \node[probebox, right=0.12cm of p3] (p4) {Probe 4\\ Cy5\\ (24 bp)};
        \node[primerbox, right=0.12cm of p4] (rev) {3' Reverse\\ Primer\\ (20 bp)};
        
        \draw[decorate, decoration={brace, amplitude=5pt, mirror}, line width=0.8pt, color=black] 
            ($(p1.south west)+(0,-0.15)$) -- ($(p4.south east)+(0,-0.15)$)
            node[midway, below=6pt, font=\small\bfseries] {Multiplex Payload (4 Probes $\times$ 24 bp)};
            
        \draw[decorate, decoration={brace, amplitude=5pt}, line width=0.8pt, color=black] 
            ($(fwd.north west)+(0,0.15)$) -- ($(rev.north east)+(0,0.15)$)
            node[midway, above=6pt, font=\small\bfseries] {Synthesized Prototype Construct (200 bp)};
    \end{tikzpicture}
    \caption{Molecular layout of the synthesized 200~bp prototype DNA taggant construct.}
    \label{fig:molecular_anatomy}
\end{figure}

\section{Physical Synthesis \& Formulation}
\begin{enumerate}
    \item \textbf{Bacterial Host Propagation}: The cloned construct was supplied in transformed \textit{Escherichia coli} (Stab Top10) on ampicillin-selective LB agar and cultured in LB broth at $37^\circ\text{C}$ for 16 hours.
    \item \textbf{Rapid Heat-Lysis Screening}: Individual bacterial colonies were suspended in $100~\mu\text{L}$ TE buffer, heat-treated at $95^\circ\text{C}$ for 10 minutes in a thermal cycler, and briefly spun. The supernatant was used directly as PCR template, eliminating the need for commercial DNA extraction kits.
    \item \textbf{Annealing Optimization \& PCR Scaling}: Gradient PCR performed across $54^\circ\text{C}\text{--}59^\circ\text{C}$ identified $57.0^\circ\text{C}$ as the optimal annealing temperature. 35 cycles of amplification produced pure 200~bp amplicon stocks.
    \item \textbf{1.0-Liter Master Spray Formulation}: 1.0~Liter of stabilized TE buffer (10~mM Tris-HCl, 1~mM EDTA, pH~8.0) was prepared with distilled water and autoclaved at $121^\circ\text{C}$ for 20 minutes. A volume of $100~\mu\text{L}$ of amplified DNA taggant (from 4--5 PCR reactions) was introduced into the 1~L carrier buffer and blended on a laboratory blood mixer to achieve complete spatial homogeneity.
\end{enumerate}

\section{Product Integration}
\begin{itemize}
    \item \textbf{Commercial Pharmaceutical QR Codes}: The taggant spray formulation was directly applied onto 2D DataMatrix and QR code panels of commercial pharmaceutical packaging (e.g., Macter syrup bottles, Bismol, and tablet cartons).
    \item \textbf{Liquid Matrices}: Directly dosed into liquid consumables (\textit{Gelsemium 30, Surkhana}) at parts-per-billion concentrations without altering color, viscosity, or chemical stability.
    \item \textbf{Industrial Roll-to-Roll Spraying}: Applied via an automated 6-station roll-to-roll spray machine with infrared flash curing ($60^\circ\text{C}$, 1.5s).
\end{itemize}

\section{Field Recovery}
\begin{enumerate}
    \item \textbf{Surface Swabbing}: A sterile swab moistened with recovery buffer wipes the printed QR code or product security zone.
    \item \textbf{Chemical Etching}: For encapsulated samples, mild fluoride buffer ($\text{NH}_4\text{F}/\text{HF}$) dissolves outer silica nano-coatings within 10 minutes.
    \item \textbf{Purification}: Standard silica spin column elutes purified DNA in $25.0~\mu\text{L}$ nuclease-free water.
\end{enumerate}

\section{qPCR Verification Protocols}
\begin{itemize}
    \item \textbf{Thermal Cycling Program}: Standard 5-stage profile: Initial denaturation at $95^\circ\text{C}$ for 5~min (1 cycle), followed by 35 cycles of denaturation ($95^\circ\text{C}$, 30s), annealing ($55^\circ\text{C}\text{--}57^\circ\text{C}$, 30s), extension ($72^\circ\text{C}$, 45s), and a final extension at $72^\circ\text{C}$ for 7~min.
    \item \textbf{4-Channel Deconvolution}: The recovered DNA is run with 4 sequence-specific fluorophore probes (FAM, HEX, Texas Red, Cy5) to simultaneously confirm the 4-probe code against the central cloud ledger.
\end{itemize}
